% !TEX root = ../main.tex

\chapter{算法框架设计与仿真环境验证}

\section{本章目标}

第三章回答的是“高层 NoMaD 策略如何设计与优化”，第四章则回答“这些候选配置如何进入 Lite3 的闭环系统”。因此，本章的重点不再是单纯的离线指标，而是统一推理模块、状态机、导航主机、中层控制映射与 MuJoCo 后端如何协同工作，并在仿真环境中形成可执行、可恢复、可扩展的闭环链路。

\section{统一闭环系统架构}

\subsection{分层设计}

结合项目实现，本文将系统划分为四层：推理层、任务组织层、中层控制层和平台执行层。推理层负责根据视觉观察与目标图像生成局部轨迹；任务组织层负责状态切换、任务队列与模式切换；中层控制层负责将轨迹映射为速度或姿态参考；平台执行层负责向 MuJoCo 或真机后端下发指令。

这种划分的关键价值在于解耦。高层策略不直接面向轮速或关节控制量，中层控制负责吸收平台差异；状态机与导航主机负责任务组织，而不直接参与模型采样细节。这样一来，同一套高层 NoMaD 推理逻辑就能被 MuJoCo 后端和真机后端复用。

\subsection{统一推理模块}

统一推理模块集中管理 \texttt{policy\_config}、\texttt{policy\_checkpoint}、\texttt{scheduler\_kind}、\texttt{ddim\_steps}、\texttt{cfg\_weight}、\texttt{tts\_budget} 等推理参数。其作用是把 DDIM、CFG、TTS 和视觉编码器替换都限制在高层策略层，使状态机、导航主机与平台后端不需要感知采样器和编码器细节。

这意味着：只要新的策略仍保持 NoMaD 兼容接口，后续系统代码就可以复用。也正因如此，本文将“推理模块”单独提炼为系统组件，而不是把 DDIM、CFG 或 TTS 的逻辑散落在多个脚本中。

\begin{figure}[htbp]
    \centering
    \fbox{\rule{0pt}{58mm}\rule{0.86\linewidth}{0pt}}
    \caption{图位占位：统一推理模块、状态机与平台后端关系图（待补）}
    \label{fig:framework-modules-placeholder}
\end{figure}

\section{状态机与导航主机设计}

\subsection{状态机职责}

在足式闭环系统中，高层推理并不是单次调用，而是需要持续嵌入控制循环。因此，系统必须回答“何时起立、何时导航、何时探索、何时恢复、何时停机”等问题。为此，本文采用显式状态机设计。当前核心状态包括 \texttt{idle}、\texttt{stand}、\texttt{navigate}、\texttt{explore}、\texttt{recovery}、\texttt{completed}、\texttt{failed} 和 \texttt{estop}。

状态机的作用不是增加流程复杂度，而是保证系统具有可恢复性和安全边界。对于 Lite3 这样的足式平台，一旦高层轨迹异常、控制映射失效或通信超时，系统必须进入恢复或停机状态，而不是继续盲目前进。

\subsection{导航主机职责}

若说状态机解决的是“任务内部如何运行”，那么导航主机解决的则是“系统当前要运行什么任务”。导航主机负责选择仿真或真机后端、装载单目标或多目标任务、切换导航/探索模式，并将任务分发给状态机执行。

这种分层使多目标导航不必由 NoMaD 一次性做全局规划，而是由导航主机把长任务拆成一系列局部子任务，再由高层策略逐段完成。由此，单目标导航、探索模式和多目标任务都能在同一高层策略之上被统一组织。

\begin{figure}[htbp]
    \centering
    \fbox{\rule{0pt}{58mm}\rule{0.86\linewidth}{0pt}}
    \caption{图位占位：状态机状态流转图（待补）}
    \label{fig:framework-fsm-placeholder}
\end{figure}

\section{MuJoCo 仿真平台设计}

\subsection{选择 MuJoCo 的原因}

MuJoCo 具有较好的接触动力学建模能力和可调仿真接口，被广泛用于足式控制与学习系统验证\cite{mujoco2012,hwangbo2019anymal,rl_locomotion_rudin2022}。对本文而言，MuJoCo 的价值不只是“把机器人动起来”，而是提供一个风险较低、接口清晰的闭环环境，用于验证高层推理、状态机、中层控制映射与平台后端是否能够协同工作。

\subsection{仿真中的系统组成}

当前 Lite3 MuJoCo 仿真链路可以概括为五部分：场景与 topomap、NoMaD 推理模块、中层轨迹映射、MuJoCo 平台接口和状态机循环。仿真后端为高层模块提供相机图像、姿态、位置信息和安全停机接口，使其尽可能与真机后端保持一致。

\subsection{仿真任务类型}

本文在 MuJoCo 中组织了四类任务：单目标导航、探索模式、状态机全流程导航和多目标任务。单目标导航用于检验基本闭环能力；探索模式用于检验 goal mask 对应的无显式目标条件在部署侧是否可用；状态机导航用于检验完整状态流是否连通；多目标任务用于检验导航主机能否持续组织任务队列。

\begin{figure}[htbp]
    \centering
    \fbox{\rule{0pt}{58mm}\rule{0.86\linewidth}{0pt}}
    \caption{图位占位：MuJoCo 仿真环境、Lite3 模型与 topomap 示例图（待补）}
    \label{fig:framework-mujoco-env-placeholder}
\end{figure}

\section{MuJoCo 闭环结果分析}

\subsection{已有结果}

当前结果表明，经过第三章筛选后的高层配置已经能够进入 MuJoCo 闭环：DDIM-2 配置可在 easy、medium、hard 地图上完成单目标导航；在 medium 场景中，适度的 CFG 能进一步改善到达效率；探索模式能够持续生成可执行轨迹；状态机任务与多目标队列也能完整运行。

\subsection{结果意义}

这些结果说明，第三章的算法层优化并没有在系统层全部失效。尤其是 DDIM-2 的闭环成功，证明扩散策略经推理加速后已经具备进入 Lite3 闭环控制区间的潜力。另一方面，多目标任务的成功表明，导航主机与状态机的分层设计是成立的，高层 NoMaD 不需要直接承担全局任务编排职责。

\subsection{结论边界}

需要强调的是，MuJoCo 结果只能证明“系统级可行性”，不能直接替代真机结论。仿真尚未完全覆盖真实相机噪声、通信时延和平台硬件扰动，因此本章的结论是：高层 NoMaD、统一推理模块、状态机、导航主机和 Lite3 平台接口已经能够构成一套稳定可运行的仿真闭环系统，并为下一章真机部署提供了可复用接口与验证基础。

\begin{figure}[htbp]
    \centering
    \fbox{\rule{0pt}{52mm}\rule{0.86\linewidth}{0pt}}
    \caption{图位占位：MuJoCo 单目标导航与多目标任务运行结果图（待补）}
    \label{fig:framework-mujoco-results-placeholder}
\end{figure}

\section{本章未完成内容与占位说明}

截至当前版本，本章尚有以下内容保留占位：MuJoCo 运行截图、状态机界面截图、失败案例分析、跨视觉编码器的正式闭环表格，以及 Tron1 相关仿真结果。这些内容不会改变本章的系统主线，但将进一步增强结果的完整性与说服力。

\section{本章小结}

本章完成了从算法层到系统层的过渡。通过统一推理模块、状态机、导航主机和 MuJoCo 后端的分层组织，NoMaD 已能够被接入 Lite3 的闭环仿真系统。MuJoCo 结果进一步说明，第三章筛选出的高层候选配置具备进入真实部署阶段的基础。下一章将在此基础上继续讨论 Orin 上的 Lite3 真机部署流程与实验组织。
